\section{Completion of the induction}\label{sec:completion}

We now assemble the two uniform statements. Their proofs have deliberately
been separated: the scalar estimates use effective Iitaka fibrations
only in smaller dimensions, whereas the denominator bound uses normal
lc indices only in smaller dimensions.

\begin{proof}[Proof of Theorem~\ref{thm:lc-index} over $\C$]
Assume $I_j,L_j$ for $j<n$. For $n=1$, a normal projective curve
with torsion canonical class is a smooth curve of genus one, so
$K_1$ holds with index one. For $n\geq2$, Proposition~\ref{prop:index-reduction}
reduces failure of $K_n$ to the terminal high-index sequence considered
in the diagonal-product argument. Proposition~\ref{prop:scalar}
supplies the required uniform section-order and Seshadri estimates.
Proposition~\ref{prop:index-contradiction} excludes that sequence.
Therefore $K_n$ holds in every dimension. Proposition~\ref{prop:lc-reduction}, together
with $L_j$ for $j<n$, now proves $L_n$ for every finite rational
coefficient set. Global ACC gives its rational DCC version as explained
in Section~\ref{sec:preliminaries}.
\end{proof}

Before passing to the Iitaka map, we spell out why no unbounded
Cartier index interferes with the comparison of section spaces.

\begin{lemma}[All-degree canonical comparison]\label{lem:all-degree}
Let $X$ be a smooth projective complex variety with pseudo-effective
$K_X$. There is a projective terminal good minimal model $V$ of $X$.
For every integer $m\geq0$, its natural birational identification of
function fields identifies
\[
          H^0(X,\OO_X(mK_X))=H^0(V,\OO_V(mK_V)).
\]
The right-hand side is a divisorial section space even when $mK_V$
is not Cartier.
\end{lemma}
\begin{proof}
Theorem~\ref{thm:good-models} and the klt MMP give a good minimal
model; starting from a smooth variety, the discrepancy comparisons
keep the model terminal. Let $W$ be a common smooth resolution. Since
$V$ is canonical, compatible canonical divisors satisfy
\[
                   K_W=q^*K_V+A,\qquad A\geq0
\]
for $q:W\to V$, with $A$ exceptional. A rational $m$-canonical form
regular as a divisorial section on $V$ has an effective zero divisor
there. This divisor is rational Cartier: it differs from $mK_V$ by
a principal divisor. Its pullback is effective, and adding $mA$
shows that the form is regular on $W$. This assertion is about actual
integral orders of the form; it does not require $mK_V$ to be Cartier.
Conversely, a form regular on $W$ is regular at the generic point of
every prime divisor on $V$, and hence is a divisorial section there
by normality. Thus all degrees agree on $V$ and $W$. The same argument
for the birational morphism $W\to X$ identifies the spaces on $W$
and $X$. Section ratios are preserved by these identifications.
\end{proof}

\begin{proof}[Proof of Theorem~\ref{thm:main} over $\C$]
We complete the simultaneous induction step. Let $X$ be a smooth
projective $n$-fold with $\kappa(X)\geq0$, and take $V$ as in
Lemma~\ref{lem:all-degree}. The semiample canonical divisor defines a
contraction $f:V\to Z$ and a rational ample divisor $D$ on $Z$ with
$K_V\sim_\Q f^*D$. This contraction represents the Iitaka fibration.

If $Z$ is a point, $K_V\sim_\Q0$, and $L_n$, just proved, gives a
common nonempty pluricanonical degree. Its image is a point.
If $\dim Z>0$, Proposition~\ref{prop:denominator} applies. Its base
adjoint is rationally linearly equivalent to $D$, so it is big.
Corollary~\ref{cor:big-base} supplies a degree depending only on $n$
whose section ratios give all of $\C(Z)$. The all-degree comparison
transfers this conclusion to $X$.

There are only finitely many possible base dimensions. Choose a common
multiple of the degrees supplied for these cases and for the point
case. Passing to this common multiple preserves the full section
field: if $s_0\ne0$ lies in the earlier degree and the new degree is
$q$ times that degree, the ratio $s/s_0$ is also the ratio
$(s s_0^{q-1})/s_0^q$ in the new complete system. All ratios in the
new system still lie in the Iitaka field by its definition. This
proves $I_n$ and closes the simultaneous induction starting from
dimension zero.
\end{proof}

\subsection{Algebraically closed fields of characteristic zero}

Let $k$ be any algebraically closed field of characteristic zero.
A given variety or pair and its finite defining data descend to a
finitely generated subfield $k_0\subset k$. Let $\overline{k_0}$ be an
algebraic closure inside $k$ and embed it into $\C$. Smoothness,
integrality, singularity conditions, and the relevant section spaces
are unchanged by these algebraically closed extensions. For
Theorem~\ref{thm:lc-index}, include a rational principalization of a
sufficiently divisible multiple among the descended data. The
uniform trivialization over $\C$ descends by the line-bundle argument
in Section~\ref{sec:preliminaries} and extends to $k$.

For Theorem~\ref{thm:main}, include one nonzero pluricanonical section
among the defining data to preserve nonnegative Kodaira dimension.
The same integer $m(n)$ gives a nonempty system after every extension.
To check its field, choose a sufficiently large divisible degree $r$
which also is divisible by $m(n)$ and defines the Iitaka fibration.
The products used in the preceding proof give a dominant rational
map from the degree-$r$ image to the degree-$m(n)$ image. It is
birational if and only if it is birational after an algebraically
closed extension: equivalently, the corresponding function-field
extension has degree one, a property detected by faithful flatness.
It is birational after comparison with $\C$, hence over $k$.
The same dimension-dependent integers therefore work over every
algebraically closed field of characteristic zero.
